\section{The generator space}

\formal{GroupPresentation.GeneratorInput; GroupPresentation.GeneratorSpace; GroupPresentation.generatorLoop}
\begin{definition}
Let $e$ denote the point of the auxiliary summand $\mathbf{1}$.  The
\emph{generator space} is $X_G=\mathrm{Susp}(\mathbf{1}\sqcup G)$, based at
its north pole.  For $g:G$, put $\ell_g=\mathrm{merid}(g)\cdot
\mathrm{merid}(e)^{-1}$.  Thus each generator determines a north-based loop.
\end{definition}

\formal{GroupPresentation.Connectedness.generatorSpaceMerePathToNorth}
\begin{lemma}
Every point of $X_G$ is merely equal to its north pole.
\end{lemma}
\begin{proof}
The proof is suspension induction.  At north it is represented by reflexivity,
and at south by the reference meridian.  Along a meridian the target is a
proposition, namely a propositional truncation of an identity type; the two
endpoint representatives therefore extend uniquely.
\end{proof}

\formal{GroupPresentation.GeneratorBouquet.generatorSpace≃∙Bouquet; GroupPresentation.GeneratorBouquet.generatorLoop-toBouquet; GroupPresentation.GeneratorBouquet.liftedGeneratorSpace≃∙Bouquet}
\begin{proposition}
There is a pointed equivalence $X_G\simeq_*\mathrm{Bouquet}(G)$ sending
$\ell_g$ to the $g$-th bouquet loop; the same assertion holds after the
universe lift used below.
\end{proposition}
\begin{proof}
The forward suspension eliminator sends both poles to the bouquet base point,
the $e$-meridian to reflexivity, and the $g$-meridian to the $g$-loop.  The
inverse sends the bouquet base point to north and its $g$-loop to $\ell_g$.
Their homotopies are respectively suspension and bouquet induction.  The
loop calculation follows by functoriality of concatenation and the unit law
for the image of the reference meridian.  The lifted proof is the same
construction with $G$ replaced by its lift.
\end{proof}
